% ====================================================================
\chapter{The Language of Mathematics}\label{ch:logic}
% ====================================================================

\begin{chapterintro}
Mathematics is written in a small, precise language. A handful of words --- \emph{proposition},
\emph{proof}, \emph{axiom}, \emph{theorem}, \emph{implies}, \emph{if and only if}, \emph{not},
\emph{for all}, \emph{there exists} --- carry almost all of the logical weight in every argument you
will meet in this module. This chapter teaches that language. Everything later in the book, from
solving an inequality to proving that a sequence has a limit, is written in it.
\end{chapterintro}

\begin{prerequisites}
None. This chapter starts from the beginning. It helps to be comfortable with elementary algebra
(rearranging an equation such as $2x+1=0$) because the examples use it.
\end{prerequisites}

\section{Propositions}\label{sec:propositions}

The basic unit of a mathematical argument is a sentence that makes a claim which is either correct or
incorrect.

\begin{definition}[Proposition][def:proposition]
A \term{proposition} is a statement that may be either true or false. \src{L1, slide 8}
\end{definition}

The definition has two parts. A proposition is a \emph{statement} --- a complete declarative sentence,
not a question, a command or a fragment --- and it must be possible to say that it is \emph{true} or
\emph{false}. It does not matter whether we know which: ``every even number greater than 2 is the sum of
two primes'' is a proposition even though nobody currently knows whether it is true.

\begin{example}[Recognising propositions]
Decide which of the following are propositions.
\begin{enumerate}
  \item $2+3=5$.
  \item $2+3=6$.
  \item $2+3$.
  \item Is $7$ greater than $4$?
  \item The unemployment rate in the United Kingdom in March 2026 was below 5\%.
  \item $x>3$.
\end{enumerate}
\solution
(a) is a proposition and it is true. (b) is a proposition and it is false; being false does not stop a
statement from being a proposition. (c) is not a statement at all: it is the name of a number, with no
verb and no claim. (d) is a question, so it is neither true nor false. (e) is a proposition: it is
either true or false as a matter of fact, whether or not you know the figure.

(f) needs care. On its own, ``$x>3$'' is neither true nor false until we know what $x$ is: it is true
when $x=5$ and false when $x=1$. Such a sentence is a \emph{proposition depending on $x$}, sometimes
called an \emph{open sentence} or a \emph{property of $x$}. It becomes a proposition once a value is
substituted for $x$, or once it is turned into a claim about many values of $x$ with a quantifier
(\cref{sec:quantifiers}).
\end{example}

\begin{note}[Propositions depending on a variable]
Much of this module works with statements such as ``$x>3$'' or ``$P_n$: $1+2+\dots+n=n(n+1)/2$'',
which depend on a variable. We write $P(x)$ or $P_n$ to show the dependence. Solving an inequality
means finding the set of all $x$ for which $P(x)$ is true (\cref{ch:ineq}); proving a statement by
induction means showing that $P_n$ is true for every natural number $n$ (\cref{ch:ind}).
\end{note}

\section{Proofs, axioms and theorems}\label{sec:proof}

\begin{definition}[Proof, axiom, theorem][def:proof]
\begin{itemize}
  \item A \term{proof} is a demonstration that one proposition follows from one or more other
  propositions.
  \item An \term{axiom} is a proposition that is accepted without proof.
  \item A \term{theorem} is a proposition that has been proved on the basis of axioms and previously
  established theorems.
\end{itemize}
\src{L1, slide 8}
\end{definition}

These three words describe how mathematical knowledge is organised. Any argument has to start
somewhere: we cannot prove everything from nothing, because every proof needs something to start from.
The starting points are the \emph{axioms}. From them we derive \emph{theorems}, by \emph{proofs} in which
each step follows logically from earlier steps. Once a theorem has been proved it may be used in later
proofs exactly as if it were an axiom.

\begin{intuition}
Think of the axioms as the rules of a game and the theorems as positions that can be reached by legal
moves. A proof is the record of the moves. Nobody asks you to justify the rules of chess; they ask you to
show that a position can be reached by playing by the rules.
\end{intuition}

In ECN115 the axioms that matter most are the properties of the real numbers listed in
\cref{ch:numbers} (labelled O1--O2, A1--A5 and M1--M6). They describe how addition, multiplication and
the order ``$>$'' behave. Everything we do with inequalities is a sequence of steps, each justified by
one of these properties. The slides make the point explicitly when they solve $2x+1>0$ by naming, at
each step, which property is being used (\cref{sec:linear-ineq}).

\begin{warning}[An axiom is not ``a fact too obvious to prove'']
An axiom is a starting assumption, chosen because we have to start somewhere. Some very obvious facts
are \emph{not} axioms in this module. For example, $0\cdot a=0$ for every real number $a$ is not on the
list of properties O1--M6; it is a theorem that can be derived from them (\cref{prop:zero-times}).
Whether something is an axiom depends on which list of assumptions has been chosen, not on how obvious
it is.
\end{warning}

\section{Implication}\label{sec:implication}

\begin{definition}[Implication][def:implication]
Given propositions $P$ and $Q$, we write $P\Rightarrow Q$ for the proposition that \term{$P$ implies $Q$},
that is, ``if $P$, then $Q$''. \src{L1, slide 8}
\end{definition}

The proposition $P$ is called the \emph{hypothesis} (or antecedent) and $Q$ the \emph{conclusion} (or
consequent). The claim $P\Rightarrow Q$ is a promise: \emph{whenever $P$ is true, $Q$ is true as well.}
It says nothing about what happens when $P$ is false.

\begin{external}[When is ``$P\Rightarrow Q$'' true?]
The slides define implication in words. In formal logic, the truth of $P\Rightarrow Q$ is fixed by the
following table, where T means true and F means false.
\begin{center}
\begin{tabular}{cc|c}
\toprule
$P$ & $Q$ & $P\Rightarrow Q$\\
\midrule
T & T & T\\
T & F & F\\
F & T & T\\
F & F & T\\
\bottomrule
\end{tabular}
\end{center}
The only way to break the promise ``if $P$ then $Q$'' is for $P$ to be true and $Q$ to be false. If $P$
is false the promise has not been tested, so the implication counts as true.
\end{external}

\begin{example}[Implications in algebra]
Consider, for real numbers $x$, the propositions $P$: ``$x>2$'' and $Q$: ``$x^2>4$''.
\begin{enumerate}
  \item Is $P\Rightarrow Q$ true for every real $x$?
  \item Is $Q\Rightarrow P$ true for every real $x$?
\end{enumerate}
\solution
(a) Yes. If $x>2$ then $x$ is positive and $x\cdot x>2\cdot x>2\cdot 2=4$ (multiplying an inequality
by a positive number keeps its direction, property M6 in \cref{ch:numbers}). So $x^2>4$.

(b) No. Take $x=-3$. Then $x^2=9>4$, so $Q$ is true, but $x>2$ is false. One such value, called a
\term{counterexample}, is enough to show that ``$Q\Rightarrow P$ for every $x$'' is false.
\end{example}

\paragraph{Ways of saying ``$P\Rightarrow Q$''} The same implication is expressed in many ways in
English, and you need to recognise all of them:
``if $P$ then $Q$'', ``$P$ implies $Q$'', ``$Q$ if $P$'', ``$Q$ follows from $P$'',
``$P$ only if $Q$'', ``$P$ is a \emph{sufficient} condition for $Q$'', and
``$Q$ is a \emph{necessary} condition for $P$''.

The words \emph{necessary} and \emph{sufficient} are used constantly in economics (for example,
``a first-order condition is necessary for an interior maximum''). In $P\Rightarrow Q$, knowing $P$ is
\emph{sufficient} to know $Q$; and $Q$ is \emph{necessary} for $P$, because $P$ cannot hold without $Q$.

\begin{warning}[``Only if'' points the same way as ``implies'']
``$P$ only if $Q$'' means $P\Rightarrow Q$, not $Q\Rightarrow P$. ``A firm makes a profit only if
price exceeds average cost'' means ``profit $\Rightarrow$ price $>$ average cost''.
\end{warning}

\subsection{Converse and contrapositive}

\begin{external}[Converse and contrapositive]
From an implication $P\Rightarrow Q$ we can build two related implications.
\begin{itemize}
  \item The \term{converse} is $Q\Rightarrow P$. It is a \emph{different} claim: the example above
  shows that $x>2\Rightarrow x^2>4$ is true while its converse is false.
  \item The \term{contrapositive} is $(\text{not }Q)\Rightarrow(\text{not }P)$. It is \emph{equivalent}
  to the original: an implication and its contrapositive are always both true or both false. For
  example, ``if $x>2$ then $x^2>4$'' has contrapositive ``if $x^2\le 4$ then $x\le 2$''.
\end{itemize}
Confusing an implication with its converse is one of the most common errors in all of mathematics.
\end{external}

\section{Equivalence}\label{sec:equivalence}

\begin{definition}[Equivalence][def:equivalence]
When $P\Rightarrow Q$ and $Q\Rightarrow P$, we write $P\Leftrightarrow Q$ and say that $P$ and $Q$ are
\term{equivalent}. \src{L1, slide 8}
\end{definition}

``$P\Leftrightarrow Q$'' is read ``$P$ if and only if $Q$'', often shortened to ``$P$ iff $Q$''. The
``if'' half is $Q\Rightarrow P$ and the ``only if'' half is $P\Rightarrow Q$. Equivalent propositions are
true in exactly the same circumstances.

\begin{core}[Why equivalence matters when solving]
When you solve an equation or inequality, you transform it step by step. If each step is an
\emph{equivalence}, the final, simple statement is true for exactly the same values of $x$ as the
original, so it describes the solution set exactly. If some step is only a one-way implication, the
final statement may be true for values that do not satisfy the original, and the ``answer'' may be wrong.
\end{core}

\begin{example}[A one-way step that introduces a false solution]
Solve $\sqrt{x+2}=x$ for real $x\ge -2$.
\solution
Squaring both sides gives $x+2=x^2$, that is, $x^2-x-2=0$, so $(x-2)(x+1)=0$ and $x=2$ or $x=-1$.
But squaring is only a one-way step: $\sqrt{x+2}=x\Rightarrow x+2=x^2$ is true, while the converse is
not (because $x$ could be negative). Checking in the original equation: $x=2$ gives $\sqrt4=2$, which
is correct; $x=-1$ gives $\sqrt1=1\ne-1$, which is wrong. The solution is $x=2$ only. The false
``solution'' $x=-1$ appeared because one step was an implication rather than an equivalence.
\end{example}

The slides illustrate the same principle with inequalities: in solving $2x+1>0$ they show not only
$2x+1>0\Rightarrow 2x>-1$ but also the reverse, $2x>-1\Rightarrow 2x+1>0$, and only then conclude
$2x+1>0\Leftrightarrow 2x>-1$ (\cref{ex:lin-ineq}).

\section{Negation}\label{sec:negation}

\begin{definition}[Negation][def:negation]
The proposition ``\term{not $P$}'' is true when the proposition $P$ is false, and ``not $P$'' is false
when $P$ is true. \src{L1, slide 8}
\end{definition}

Negation reverses the truth value. Forming the negation of a mathematical statement correctly is a skill
in its own right, needed for proofs by contradiction and for showing that something does \emph{not}
happen (for example, that a sequence does not have a limit, \cref{ex:no-limit}).

\paragraph{Negating an inequality} For real numbers, exactly one of $a>b$, $a=b$ and $a<b$ is true
(property O1, \cref{ch:numbers}). Therefore
\[
\text{not}(x>0)\quad\Longleftrightarrow\quad x\le 0, \qquad\qquad \text{not}(x\le 3)\quad\Longleftrightarrow\quad x>3.
\]
The negation of ``$x>0$'' is \emph{not} ``$x<0$'': it must also include $x=0$.

\begin{external}[Negating ``and'' and ``or'': De Morgan's laws]
The study guide adds the two rules for negating compound statements, known as De Morgan's laws:
\[
\text{not}(P\text{ and }Q)\ \Leftrightarrow\ (\text{not }P)\text{ or }(\text{not }Q),\qquad
\text{not}(P\text{ or }Q)\ \Leftrightarrow\ (\text{not }P)\text{ and }(\text{not }Q).
\]
For example, the negation of ``$x>1$ and $y>1$'' is ``$x\le1$ or $y\le1$''. In mathematics ``or'' is
always \emph{inclusive}: ``$P$ or $Q$'' is true when at least one of them is true, including when both are.
Finally, the negation of an implication is
$\text{not}(P\Rightarrow Q)\Leftrightarrow P\text{ and }(\text{not }Q)$: to deny a promise you must show
that the condition held and the promised outcome did not.
\end{external}

\section{Quantifiers: ``for all'' and ``there exists''}\label{sec:quantifiers}

Statements about \emph{many} objects at once need two further symbols.

\begin{definition}[Quantifiers][def:quantifiers]
Let $A$ be a set and $P$ a property (a proposition depending on an element $x$).
\begin{itemize}
  \item ``$\exists x\in A: P$'', or ``$\exists x\in A$ such that $P$'', means \emph{there exists at least
  one element $x$ in the set $A$ such that the property $P$ holds for $x$}. The symbol $\exists$ is the
  \term{existential quantifier}.
  \item ``$\forall x\in A,\ P$'' means \emph{for all elements $x$ of the set $A$, the property $P$ is
  true}. The symbol $\forall$ is the \term{universal quantifier}.
\end{itemize}
\src{L2, slide 21}
\end{definition}

\begin{example}[Quantifiers on a finite set]\label{ex:quant-finite}
Let $A=\{2,4,6,8,10,12\}$. Show that $\exists x\in A: x\ge 10$ and that $\forall x\in A,\ x<20$.
\solution
For the first statement it is enough to exhibit one element: $10\in A$ and $10\ge10$ (so is $12$). For
the second we must check every element: $2,4,6,8,10,12$ are all less than $20$. Both statements are
therefore true. \src{L2, slide 21}
\end{example}

\begin{core}[How to prove and disprove quantified statements]
\begin{itemize}
  \item To prove $\exists x\in A: P(x)$, exhibit one particular $x\in A$ and check that $P(x)$ holds.
  \item To prove $\forall x\in A,\ P(x)$, take an \emph{arbitrary} $x\in A$ --- one about which you assume
  nothing except that it belongs to $A$ --- and show that $P(x)$ holds. Checking examples is not enough
  unless $A$ is finite and you check every element.
  \item To disprove $\forall x\in A,\ P(x)$, find one counterexample: an $x\in A$ for which $P(x)$ is
  false.
  \item To disprove $\exists x\in A: P(x)$, show that $P(x)$ is false for every $x\in A$.
\end{itemize}
\end{core}

\subsection{Negating quantified statements}

\begin{external}[Negation swaps the quantifiers]
\[
\text{not}\big(\forall x\in A,\ P(x)\big)\ \Leftrightarrow\ \exists x\in A:\ \text{not }P(x),\qquad
\text{not}\big(\exists x\in A: P(x)\big)\ \Leftrightarrow\ \forall x\in A,\ \text{not }P(x).
\]
``Not every student passed'' means ``some student did not pass''; ``there is no student who passed''
means ``every student failed''. When negating, the set $A$ after the quantifier does not change: the
negation of a claim about real numbers is still a claim about real numbers.
\end{external}

\subsection{The order of quantifiers}

When a statement contains several quantifiers, their order matters.

\begin{example}[Order of quantifiers]\label{ex:quant-order}
Compare
\[
\text{(i)}\ \ \forall n\in\Z,\ \exists m\in\Z:\ m>n
\qquad\text{and}\qquad
\text{(ii)}\ \ \exists m\in\Z:\ \forall n\in\Z,\ m>n.
\]
\solution
Statement (i) says that for every integer $n$ there is some integer larger than it. This is true: given
$n$, take $m=n+1$. The choice of $m$ is allowed to depend on $n$.

Statement (ii) says that there is a single integer $m$ that is larger than every integer. This is false:
whatever $m$ is proposed, the integer $n=m$ is not smaller than it. Here $m$ must be chosen first and
then work for every $n$ simultaneously.
\end{example}

The definition of the limit of a sequence (\cref{ch:limits}) has the shape
``$\forall\varepsilon>0,\ \exists N\in\N:\ \forall n>N,\ \dots$''. Reading it correctly depends entirely
on this point: $N$ comes after $\varepsilon$, so $N$ may depend on $\varepsilon$.

\section{How proofs are built}\label{sec:proof-types}

The proofs in this book use a small number of patterns. They are introduced in the chapters where they
are first needed, and are listed here so that you can recognise them.

\begin{description}[style=nextline, font=\sffamily\bfseries\color{accent}]
\item[Direct proof (chain of implications)] Start from what is given and reach the conclusion by steps
each of which follows from earlier ones. Used for almost every computation, for example solving an
inequality (\cref{ch:ineq}).
\item[Proof of an equivalence] Prove $P\Rightarrow Q$ and $Q\Rightarrow P$ separately, or build a chain
of steps each of which is reversible.
\item[Proof by cases] If one of several cases must hold, and the conclusion follows in each case, the
conclusion holds. Used for the sign analysis of a product (\cref{ex:product-ineq}) and for absolute
values (\cref{ch:abs}).
\item[Proof by contradiction] To prove $P$, assume ``not $P$'' and derive something impossible (for
example $1<\tfrac12$). Then ``not $P$'' cannot hold, so $P$ does. Used to show that the sequence
$1,0,1,0,\dots$ has no limit (\cref{ex:no-limit}).
\item[Proof by mathematical induction] To prove $P_n$ for every $n\in\N$, prove $P_1$ and prove that
$P_n\Rightarrow P_{n+1}$ for every $n$ (\cref{ch:ind}).
\end{description}

% --------------------------------------------------------------------
\section{Chapter review}
% --------------------------------------------------------------------
\subsection*{Summary}
A proposition is a statement that is either true or false. Mathematics organises propositions into
axioms (accepted without proof) and theorems (proved from axioms and earlier theorems), connected by
proofs. Propositions are combined with ``implies'' ($\Rightarrow$), ``if and only if'' ($\Leftrightarrow$)
and ``not''. An implication is a one-way promise; an equivalence is a two-way one, and only a chain of
equivalences is guaranteed to preserve a solution set exactly. Quantifiers turn properties of an element
into claims about whole sets: $\exists$ is proved by an example, $\forall$ by an argument about an
arbitrary element, and a false $\forall$ is disproved by one counterexample. Negation swaps $\forall$ and
$\exists$, and the order of mixed quantifiers changes the meaning.

\subsection*{Key definitions}
\begin{itemize}
  \item \textbf{Proposition:} a statement that may be either true or false.
  \item \textbf{Proof:} a demonstration that one proposition follows from one or more others.
  \item \textbf{Axiom:} a proposition accepted without proof.
  \item \textbf{Theorem:} a proposition proved from axioms and previously established theorems.
  \item \textbf{$P\Rightarrow Q$:} if $P$ then $Q$.\quad \textbf{$P\Leftrightarrow Q$:} $P\Rightarrow Q$ and
  $Q\Rightarrow P$.\quad \textbf{not $P$:} true exactly when $P$ is false.
  \item \textbf{$\exists x\in A: P$:} at least one $x\in A$ has property $P$.\quad
  \textbf{$\forall x\in A,\ P$:} every $x\in A$ has property $P$.
\end{itemize}

\subsection*{Key logical equivalences}
\begin{gather*}
P\Leftrightarrow Q\ \text{ means }\ (P\Rightarrow Q)\text{ and }(Q\Rightarrow P);\qquad
\text{not}(x>a)\Leftrightarrow x\le a;\\
\text{not}(\forall x\in A,\ P)\Leftrightarrow \exists x\in A:\ \text{not }P;\qquad
\text{not}(\exists x\in A: P)\Leftrightarrow \forall x\in A,\ \text{not }P.
\end{gather*}

\subsection*{Connections}
The vocabulary of this chapter is used throughout. Set equality is defined by a two-way condition
(\cref{ch:sets}); the properties of the real numbers are axioms (\cref{ch:numbers}); solving
inequalities is the construction of a chain of equivalences (\cref{ch:ineq}); induction is a way of
proving a $\forall$ statement about $\N$ (\cref{ch:ind}); and the definition of a limit is a statement
with three nested quantifiers (\cref{ch:limits}).

\subsection*{Common mistakes}
\begin{itemize}
  \item Treating $P\Rightarrow Q$ as if it meant $Q\Rightarrow P$ (confusing an implication with its converse).
  \item Reading ``$P$ only if $Q$'' as $Q\Rightarrow P$.
  \item Negating $x>0$ as $x<0$ instead of $x\le 0$.
  \item ``Proving'' a $\forall$ statement by checking a few examples.
  \item Swapping the order of $\forall$ and $\exists$.
  \item Using a one-way step (such as squaring) when solving and not checking the answers.
\end{itemize}

\subsection*{Exam checklist}
\begin{checklist}
  \item Define proposition, proof, axiom and theorem in the words of the course.
  \item Explain the difference between $\Rightarrow$ and $\Leftrightarrow$, with an example.
  \item Form the negation of an inequality and of a quantified statement.
  \item Prove an $\exists$ statement with an example and disprove a $\forall$ statement with a counterexample.
  \item Explain why the order of quantifiers matters, using \cref{ex:quant-order}.
\end{checklist}

% --------------------------------------------------------------------
\section{Practice questions}
% --------------------------------------------------------------------
\pq[basic] Which of the following are propositions? For each proposition, say whether it is true or false.
(a) $7$ is an even number. (b) $3x+2$. (c) $\sqrt{16}=4$. (d) Calculate $5\times 6$.
(e) Every natural number is greater than $0$.

\pq[basic] State the converse and the contrapositive of: ``If a firm's revenue exceeds its costs, then the
firm makes a profit.''

\pq[conceptual] Explain why checking that a statement holds for $n=1,2,\dots,1000$ does not prove that it
holds for every natural number $n$.

\pq[mathematical] Write the negation of each statement, simplifying as far as possible.
(a) $x\ge 5$. (b) $\forall x\in\R,\ x^2\ge 0$. (c) $\exists n\in\N: n^2=2$.
(d) $x>1$ and $x<4$.

\pq[mathematical] Let $A=\{1,3,5,7,9\}$. Decide whether each statement is true, with a reason.
(a) $\forall x\in A,\ x$ is odd. (b) $\exists x\in A: x>8$. (c) $\forall x\in A,\ x<9$.
(d) $\exists x\in A: x^2=25$.

\pq[application] An economist claims: ``For every price $p>0$ there is a quantity $q>0$ that the consumer
is willing to buy.'' A second economist claims: ``There is a quantity $q>0$ that the consumer is willing to
buy at every price $p>0$.'' Write both claims with quantifiers and explain why they say different things.

\pq[multi-step] Consider the claim ``for all real $x$, if $x^2=9$ then $x=3$''. (a) Is it true? Justify.
(b) State its converse and decide whether the converse is true. (c) Replace the conclusion so that the
resulting implication is true and its converse is also true, and write it as an equivalence.

\pq[exam-style] (a) Explain what is meant by a \emph{proposition}, an \emph{axiom} and a \emph{theorem}.
(b) Let $A=\{x\in\R: 0<x<1\}$. Determine whether each of the following is true, giving a proof or a
counterexample: (i) $\forall x\in A,\ x^2<x$; (ii) $\exists x\in A: x^2=x$.

% --------------------------------------------------------------------
\section{Solutions to practice questions}
% --------------------------------------------------------------------
\pqsol{1.1}
(a) A proposition; false. (b) Not a proposition: it is an expression (a number once $x$ is known), not a
statement. (c) A proposition; true. (d) Not a proposition: it is an instruction. (e) A proposition; true
($\N=\{1,2,3,\dots\}$, \cref{ch:numbers}).

\pqsol{1.2}
Let $P$ be ``revenue exceeds costs'' and $Q$ be ``the firm makes a profit''. The converse is ``If the firm
makes a profit, then its revenue exceeds its costs.'' The contrapositive is ``If the firm does not make a
profit, then its revenue does not exceed its costs.'' (With profit defined as revenue minus costs, all
three statements are true here, but in general only the contrapositive is guaranteed to have the same
truth value as the original.)

\pqsol{1.3}
The statement is a claim about infinitely many propositions $P_1,P_2,P_3,\dots$. Checking $1000$ of them
establishes those $1000$ and says nothing about $P_{1001}$ or any later one; however many cases are
checked, infinitely many remain. A proof must cover an \emph{arbitrary} natural number, for example by
mathematical induction (\cref{ch:ind}). There are famous statements that hold for very many initial cases
and then fail.

\pqsol{1.4}
(a) $x<5$. (b) $\exists x\in\R: x^2<0$ (a false statement, as it should be, since the original is true).
(c) $\forall n\in\N,\ n^2\ne 2$. (d) By De Morgan's law, ``$x\le 1$ or $x\ge 4$''.

\pqsol{1.5}
(a) True: $1,3,5,7,9$ are all odd. (b) True: $9\in A$ and $9>8$. (c) False: $9\in A$ but $9<9$ is false;
$x=9$ is a counterexample. (d) True: $x=5$.

\pqsol{1.6}
The first claim is $\forall p>0,\ \exists q>0$: the consumer buys $q$ at price $p$. The quantity may depend
on the price, which is how demand normally works. The second is $\exists q>0:\ \forall p>0$, the consumer
buys $q$ at price $p$: one fixed quantity would be bought however high the price, which is a far stronger
(and economically implausible) claim. The quantifiers are the same; only their order differs, and that
changes the meaning (\cref{ex:quant-order}).

\pqsol{1.7}
(a) False: $x=-3$ satisfies $x^2=9$ but $x\ne3$. (b) The converse is ``if $x=3$ then $x^2=9$'', which is
true. (c) ``If $x^2=9$ then $x=3$ or $x=-3$'' is true (since $x^2-9=(x-3)(x+3)=0$ forces one factor to be
zero), and its converse is true as well, so
$x^2=9\Leftrightarrow (x=3\text{ or }x=-3)$.

\pqsol{1.8}
(a) A proposition is a statement that may be either true or false. An axiom is a proposition accepted
without proof. A theorem is a proposition that has been proved on the basis of axioms and previously
established theorems.
(b)(i) True. Take an arbitrary $x\in A$, so $0<x<1$. Multiplying $x<1$ by the positive number $x$ keeps the
direction (M6): $x\cdot x<1\cdot x$, that is, $x^2<x$. Since $x$ was arbitrary, the statement holds for all
$x\in A$.
(ii) False. If $x^2=x$ then $x(x-1)=0$, so $x=0$ or $x=1$; neither lies in $A$. Hence no $x\in A$ satisfies
$x^2=x$. (Alternatively, (ii) contradicts (i).)
